% mass-gap-discussion.tex — Chapter on the sharpness, physical-unit
% calibration, lattice-glueball comparison, and the SU(N) extension
% framework that ends Part V.

\chapter{Sharpness, Comparison, and the
  \texorpdfstring{$\SUN$}{SU(N)} Extension}
\label{ch:mass-gap-discussion}

\noindent\emph{The constructive bound $\massgap \ge 999/1000$ in
context: comparison with lattice-Monte-Carlo glueball spectroscopy,
physical-unit calibration, sharpness discussion, and the analytic
roadmap to $\SUthree$.}

This chapter places the headline theorem of
\Cref{ch:mass-gap-theorem} ($\sigma(\Hop) \cap (0, \massgap) =
\emptyset$ with $\massgap \ge 999/1000$ in dimensionless units) in
the physical context: what the bound says about the actual
$\SUtwo$ mass gap, what it does not say, and how the framework
generalizes to $\SUthree$ and beyond.  The discussion is
intentionally conservative: the constructive bound is positive and
$\Qp$-rational, but it is several orders of magnitude looser than
the empirical glueball mass; the gap-vs-tightness story is
unpacked here.

\section{The bound in context}
\label{sec:headline-in-context}

\marginkey{Strict, not sharp}%

The bound $\massgap \ge 999/1000$ is in dimensionless units where
the lattice-scale-$0$ string tension $\sigma_0$ is normalized to $1$.
What this means physically:

\begin{itemize}
  \item \textbf{Positive.}  $\massgap > 0$ is the substantive
    content of the Clay statement; the constructive bound is
    strictly positive at every operating point inside the
    convergent regime.
  \item \textbf{$\Qp$-rational.}  Every constant in the manuscript
    is a rational number with a small denominator; the bound is
    not a symbolic limit of irrational quantities but a concrete
    fraction.  The Lean witness verifies the bound to the literal
    rational $999/1000$.
  \item \textbf{Lattice-scale-zero normalized.}  The number $999/1000$
    is the dimensionless ratio $\massgap / \sqrt{\sigma_0}$; the
    physical mass in MeV depends on the calibration of
    $\sqrt{\sigma_0}$, discussed in
    \Cref{sec:physical-units} below.
\end{itemize}

\section{Comparison with lattice glueball spectroscopy}
\label{sec:teper-comparison}

\marginkey{Empirical $0^{++} / \sqrt{\sigma} \approx 3.74$}%

Lattice Monte Carlo spectroscopy has produced detailed measurements
of the $\SUtwo$ glueball spectrum.  The lightest glueball, the
scalar $0^{++}$, sets the empirical mass gap.

\begin{theorem}[Empirical $\SUtwo$ glueball-to-string-tension ratio,
  Teper et~al.]
\label{thm:teper-ratio}
Continuum-extrapolated lattice Monte Carlo measurements of the
$\SUtwo$ scalar glueball mass give
\begin{equation}
  \label{eq:teper-ratio}
  \frac{m(0^{++})}{\sqrt{\sigma}} \;\approx\; 3.74
  \quad \mathrm{\cite{Teper1998Glueballs}},
  \qquad
  \frac{m(0^{++})}{\sqrt{\sigma}} \;\approx\; 3.78 \pm 0.07
  \quad \mathrm{\cite{LuciniTeperWenger2004}} ,
\end{equation}
where $\sigma$ is the continuum string tension.
\end{theorem}

Compared to the manuscript bound $\massgap \ge 999/1000 \approx 1.0$:

\begin{infobox}[main][title=Bound versus empirical $\SUtwo$ mass]
\centering
\renewcommand{\arraystretch}{1.25}
\begin{tabular}{lc}
\toprule
Quantity                                         & Value     \\
\midrule
This manuscript's constructive bound $\massgap$  & $\ge 999/1000 \approx 1.0$ \\
Empirical $m(0^{++}) / \sqrt{\sigma}$ (Teper)    & $\approx 3.74$ \\
Looseness factor                                 & $\approx 3.74$ \\
\bottomrule
\end{tabular}
\end{infobox}

The constructive bound is loose by a factor of $\sim 3.7$ relative
to the empirical glueball mass.  This looseness is expected: the
chain $\Qp\text{-data} \Rightarrow \sighinf \Rightarrow \mathrm{OS4}
\Rightarrow \massgap$ uses the OS4 cluster rate $\sighinf$ as the
spectral lower bound, which is the linear-decay-to-gap relation
$m \ge \sqrt{\sighinf}$ rather than a sharper estimate.  Removing
the factor of $\sim 3.7$ would require a Brydges--Federbush-style
sharper-decay-to-gap argument; an open analytical refinement.

\section{Physical-unit calibration}
\label{sec:physical-units}

\marginkey{Pure-gauge $\sqrt\sigma \approx 440$ MeV}%

To convert the dimensionless $\massgap$ to MeV, fix the continuum
string tension at its standard pure-gauge calibration
\begin{equation}
  \label{eq:string-calibration}
  \sqrt{\sigma} \;\approx\; 440 \text{ MeV}
  \quad \text{(pure-gauge, lattice-standard convention)}.
\end{equation}
Under this calibration:

\begin{corollary}[$\massgap$ in MeV]
\label{cor:m0-in-mev}
The constructive bound translates to
\begin{equation}
  \label{eq:m0-in-mev}
  \massgap \cdot \sqrt{\sigma}
  \;\ge\;
  (999/1000) \cdot 440 \text{ MeV}
  \;\approx\;
  439.6 \text{ MeV},
\end{equation}
which is several orders of magnitude below the empirical glueball
mass of $\approx 1.5{-}1.6$ GeV but is, importantly, strictly
positive.
\end{corollary}

\section{Where the looseness comes from}
\label{sec:looseness-sources}

The factor-$3.7$ gap between the constructive bound and the
empirical mass decomposes into named contributions, each a
potential refinement target.

\begin{enumerate}[label=(\arabic*)]
  \item \textbf{Linear-decay-to-gap step.}
    The OS4 cluster rate of \Cref{ch:os-axioms} is identified with
    the spectral gap via $m \ge m_0^{(\mathrm{OS4})}$ (\Cref{ch:os-reconstruction}'s
    cluster-to-gap bridge), but the actual spectral gap can be
    significantly larger than the cluster rate when the propagator
    has additional structure.  Estimated tightening:
    factor $\sim 2$ if a Brydges--Federbush-style sharper
    decay-to-gap relation were established.
  \item \textbf{$\Kcap$ union-bound looseness.}
    \Cref{ch:overlap-kp},
    Conjecture~\ref{conj:K-cap-sharpness}, observed that
    $\Kcap \le \Ksf + \Klf$ is loose by a factor up to $\approx
    1.18$ on the canonical trajectory.  A sharper
    sheaf-domination-gluing bound would tighten the chain at this
    step.
  \item \textbf{Polymer-residual $\cpolymer$ overcounting.}
    \Cref{ch:sigma-hat-trajectory}'s derived
    $\cpolymer = 2$ is the worst-case channel multiplicity from the
    triangle inequality.  A sign-preservation lemma on the
    cluster-expansion residual would reduce $\cpolymer$ to $1$,
    halving the per-step $\eta$ and improving $\massgap$ to
    $\sqrt{1 - 9/(2 \cdot 10000)} \approx 0.99955$ (relative to
    $0.999$).
  \item \textbf{Fiducial-bound looseness.}
    The named fiducials $\clambda$, $\muplaq$, $\cpolymer$, $\Amin$
    each contribute an over-estimation factor; the joint
    sensitivity is bounded by the bootstrap analysis of
    Appendix~\ref{ch:appx-fiducials}, but no individual fiducial
    refinement closes the factor-$3.7$ gap.
\end{enumerate}

Each tightening target is an open analytical refinement, not a
blocker for the manuscript's existence claim.

\section{The $\SUN$ extension framework}
\label{sec:su-n-extension}

\marginkey{Roadmap to $\SUthree$}%

The manuscript framework is group-agnostic at the abstract level:
the inductive state $\Ical_k$, the block-spin morphism $\Tb$, the
$\sigh$-budget closure, and the OS-reconstruction chain all hold
for any compact simple gauge group $G$.  Specialization to a
specific $G$ enters at three named points:

\begin{enumerate}[label=(E\arabic*)]
  \item \textbf{Representation theory.}
    \Cref{ch:su2-rep-theory} fixes the irrep indexing,
    dimensions, Casimirs.  For $\SUN$ the representations are
    indexed by Young tableaux (or highest weights), and the
    smallest nonzero quadratic Casimir is
    $\Casimir(\mathrm{fundamental}) = (N^2 - 1)/(2N)$.
    For $\SUthree$: $\Casimir(\mathrm{fundamental}) = 4/3$.
  \item \textbf{Heat-kernel character expansion.}
    \Cref{ch:heat-kernel}'s character-expansion structure is
    universal across compact simple groups; the SU(2)-specific
    closed-form character $\chi_j(\theta) = \sin((2j+1)\theta/2)
    / \sin(\theta/2)$ is replaced by a multi-variable Weyl
    character formula at higher rank.
  \item \textbf{Plaquette-polymer connectivity $\muplaq$.}
    The plaquette graph structure depends only on the lattice
    dimension $d$, not on the gauge group, so $\muplaq$ of
    \Cref{ch:large-field-kp},
    Convention~\ref{conv:mu-plaq}, is unchanged.
\end{enumerate}

The structural part of the manuscript (Parts III--V) carries over
to $\SUthree$ with no architectural change.  What changes is the
\emph{numerical} pilot: the smaller Casimir gap of $\SUthree$ —
$\Casimir(\mathrm{fundamental}) = 4/3$ versus the $\SUtwo$ value of
$\tfrac34$ — gives a different convergent-regime threshold
$g_{\max}^{\mathrm{sf}}$ and a different aggregate $\eta$ budget.
\emph{Preliminary expectation}: a $\SUthree$ pilot at appropriate
$g_0$ would close the inductive step with a different but still
$\Qp$-rational bound $\massgap_{\SUthree}$.

\begin{conjecture}[$\SUthree$ pilot closure]
\label{conj:su3-pilot}
There exists $g_0 \in (0, 1)$ such that the $\SUthree$ pilot
validator (the natural analog of \Cref{ch:pilot-su2},
Definition~\ref{def:pilot-validator}) closes at $K = 3$,
yielding a $\Qp$-rational $\massgap_{\SUthree} > 0$ via the same
inductive-step-plus-continuum-limit-plus-OS-reconstruction chain.
\end{conjecture}

The conjecture is reachable with the current framework; the work
required is mechanical (rerun the Phase-1--3 numerics with the
$\SUthree$ representation-theoretic and heat-kernel inputs) and
should not introduce any new analytic open question.

\section{What this manuscript does not deliver}
\label{sec:non-deliverables}

\begin{itemize}
  \item \textbf{A sharp mass-gap value.}  The constructive bound is
    a strict lower bound, several orders of magnitude looser than
    the empirical mass.  A sharp constructive value of $\massgap$
    requires the analytical refinements of
    \Cref{sec:looseness-sources}.
  \item \textbf{Confinement.}  The Clay-statement bonus condition
    (C-bis) on $\SUthree$ confinement is not addressed by this
    manuscript.  Confinement requires showing area-law behavior of
    Wilson loops at \emph{all} sizes, not just in the small-field
    regime; the framework here delivers area-law in the
    convergent regime but does not extend to the strong-coupling
    side.
  \item \textbf{Glueball-spectrum identification.}  The bound
    $\massgap$ is the spectral gap of $\Hop$, not the
    lightest-glueball mass per se.  The
    spin-and-parity-resolved spectroscopy is beyond the present
    scope; bridging from spectral-gap-of-$\Hop$ to specific
    glueball quantum numbers requires additional analytic input.
  \item \textbf{Rotation-invariance restoration in OS1.}
    \Cref{ch:os-axioms} addressed only the hypercubic $H_4$
    subgroup of $O(4)$; the full continuum rotation invariance
    requires a separate restoration argument, deferred to future
    work.
\end{itemize}

\section{What it does deliver}
\label{sec:deliverables}

The manuscript closes the loop on the $\SUtwo$ instance of
(C1)--(C3) of \Cref{ch:introduction},
\Cref{sec:clay-statement}:

\begin{itemize}
  \item \textbf{Existence (C1).}  The projective-limit construction
    of $\muinf$ on $\Ominf$
    (\Cref{ch:continuum-measure}) delivers a non-Gaussian
    probability measure on $\SUtwo$-gauge fields over $\mathbb{R}^4$;
    OS-axiom verification (\Cref{ch:os-axioms}) and
    OS-reconstruction (\Cref{ch:os-reconstruction}) give the
    associated quantum field theory.
  \item \textbf{Non-triviality (C2).}  The constructed $\muinf$
    is non-Gaussian by construction (the heat-kernel Wilson
    measure with explicit positive coupling $g_k > 0$).
  \item \textbf{Mass gap (C3).}  $\sigma(\Hop) \cap (0, \massgap)
    = \emptyset$ with $\massgap \ge 999/1000$ in dimensionless
    units, certified by a Lean witness with
    axiom-gate $\codepath{[propext]}$ alone
    (\Cref{ch:mass-gap-theorem}).
\end{itemize}

That is the Clay-statement headline, specialized to $\SUtwo$, with
every constant pinned to an explicit $\Qp$-rational value and every
analytical input tracked through the manuscript's code provenance
annotations.

\section{Outlook}
\label{sec:outlook}

The natural next directions are:
\begin{itemize}
  \item Close the $\SUthree$ pilot
    (\cref{conj:su3-pilot}) by rerunning the numerical
    machinery with the appropriate representation-theoretic inputs.
  \item Tighten the $\massgap$ bound by addressing the
    looseness sources of \Cref{sec:looseness-sources}, most
    productively (3)~the polymer-residual $\cpolymer$ refinement
    and~(1) the sharper-decay-to-gap relation.
  \item Build the constructive-reals layer that would extend the
    $\Qp$-rational bounds of the manuscript to a real-valued
    continuum $\massgap$ — closing the gap between the
    finite-truncation rational bound and the asymptotic real-valued
    continuum-limit value.
  \item Resolve the full $O(4)$ rotation-invariance restoration of
    \Cref{ch:os-axioms},
    \Cref{sec:rotation-restoration}, currently open.
\end{itemize}

Each is a well-defined refinement task on top of the framework laid
out here.  None changes the existence statement: a constructive,
strictly positive, $\Qp$-rational $\SUtwo$ mass gap has been
established.
